\documentclass[11pt]{article}
\usepackage[utf8]{inputenc}
\usepackage{amsmath,amssymb,amsthm}
\usepackage[margin=1in]{geometry}
\usepackage{microtype}
\usepackage{hyperref}
\usepackage{xcolor}
\usepackage{enumitem}
\usepackage{newunicodechar}
\emergencystretch=3em
\newunicodechar{ℝ}{\ensuremath{\mathbb{R}}}
\newunicodechar{λ}{\ensuremath{\lambda}}
\newunicodechar{κ}{\ensuremath{\kappa}}
\newunicodechar{·}{\ensuremath{\cdot}}
\newunicodechar{∂}{\ensuremath{\partial}}

\newcommand{\deriv}{\operatorname{deriv}}
\newcommand{\tideref}[2][]{\href{https://therisensea.org/tides/#2/}{\ifx&#1&\texttt{#2}\else#1\fi}}

\title{Tide retrospective:\\\emph{cumulants as derivatives of the log-partition}}
\author{Claude Opus 4.8 (1M ctx), with Daniel Murfet}
\date{2026-05-30}

\begin{document}
\maketitle

\begin{abstract}
The synthesis that closes the anharmonic cumulant arc: the connected cumulants
$\kappa_2,\kappa_3,\kappa_4$ --- proved across the arc as derivatives of the
perturbed \emph{mean} $M=G_1/G_0$ --- are identified as derivatives of the
genuine cumulant generating function $K=\log Z$ (with $Z=G_0$):
\begin{align*}
  &\mathrm{iteratedDeriv}\,(n{+}1)\,K\,0 = \mathrm{iteratedDeriv}\,n\,M\,0, \\
  &\quad\text{so}\quad
  \kappa_2 = \partial^2 K|_0,\quad \kappa_3 = \partial^3 K|_0,\quad
  \kappa_4 = \partial^4 K|_0 .
\end{align*}
The new file\footnote{\texttt{AnharmonicCumulantsLogPartition.lean}.} builds
clean first try, no \texttt{sorry}/\texttt{axiom}/\texttt{native\_decide}.
\end{abstract}

\section{Setting}

The tenth and final tide of a single-day anharmonic arc on the laplace seabed.
Across the arc the cumulants were computed as derivatives of the mean
$M(h)=\langle u\rangle_h=G_1(h)/G_0(h)$ (because $K'=M$, the
fluctuation--dissipation route is the most direct). This tide records the
``why'': $K=\log Z$ is the cumulant \emph{generating} function, so its
derivatives at $0$ \emph{are} the connected cumulants. Closing this loop --- by
relating $\partial^{n+1}\log Z$ to $\partial^n M$ --- restores the textbook
definition and ties the whole arc back to its first rung,
\texttt{logPartition\_hasDerivAt}.

\section{The thing formalised}

Writing $K(h)=\log Z(h)=\log G_0(h)$ and $M(h)=G_1(h)/G_0(h)$:
\begin{quote}\itshape
For the anharmonic potential, $\lambda,\gamma>0$, $\alpha^2<3\lambda\gamma$,
$t>0$: $\ \mathrm{iteratedDeriv}\,(n{+}1)\,K\,0 = \mathrm{iteratedDeriv}\,n\,M\,0$
for all $n$; and the named corollaries
$\partial^2 K|_0=\kappa_2$, $\partial^3 K|_0=\kappa_3$,
$\partial^4 K|_0=\kappa_4$ (each in the explicit $G_0,\dots,G_4$ closed form of
the corresponding cumulant tide).
\end{quote}

\section{Proof strategy}

Three short steps. \emph{(i)} The general-$h$ first derivative of $K$ is the
mean: \texttt{HasDerivAt.log} on the partition derivative
$\mathrm{HasDerivAt}\,G_0\,(G_1(h_0))\,h_0$ with $G_0(h_0)>0$ gives
$K'(h_0)=G_1(h_0)/G_0(h_0)$, hence $\deriv K =_{\text{ev}} M$ on
$\mathrm{ball}\,0\,1$. \emph{(ii)} The ladder lemma: by \texttt{iteratedDeriv\_succ'},
$\mathrm{iteratedDeriv}\,(n{+}1)\,K = \mathrm{iteratedDeriv}\,n\,(\deriv K)$, and
\texttt{Filter.EventuallyEq.iteratedDeriv\_eq} rewrites $\deriv K$ to $M$, giving
$\mathrm{iteratedDeriv}\,(n{+}1)\,K\,0 = \mathrm{iteratedDeriv}\,n\,M\,0$.
\emph{(iii)} The three corollaries are \texttt{rw [ladder]} then \texttt{exact}
the existing cumulant theorems ($\kappa_2$ at $n=1$ via \texttt{iteratedDeriv\_one},
$\kappa_3$ at $n=2$, $\kappa_4$ at $n=3$).

\section{Roadblocks and resolutions}

None --- first-try build. The synthesis is pure chaining of established
lemmas, and crucially \texttt{Filter.EventuallyEq.iteratedDeriv\_eq} already
exists in Mathlib, so the ``derivatives respect germs'' step needed no bespoke
induction. Because the corollaries discharge by \texttt{exact}-ing the cumulant
theorems, Lean checks that the stated closed forms match the proved ones ---
no risk of a silently mistyped $\kappa_4$.

\section{What was Mathlib, what was new}

\textbf{Mathlib.} \texttt{HasDerivAt.log}, \texttt{HasDerivAt.deriv};
\texttt{iteratedDeriv\_succ'}, \texttt{iteratedDeriv\_one};
\texttt{Filter.EventuallyEq.iteratedDeriv\_eq}.

\textbf{Seabed (reused).} The partition derivative and positivity, and the
three cumulant theorems --- all from this session's arc.

\textbf{New.} The general-$h$ log-partition derivative, the ladder lemma, and
the three $\partial^n\log Z$ corollaries. About 130 lines.

\section{Lessons}

\begin{itemize}[itemsep=3pt]
\item \textbf{Two routes to the cumulants, one cheap bridge.} Computing
  $\kappa_n$ as $\partial^{n-1} M$ (fluctuation--dissipation) sidesteps the
  $\log$ entirely and was the right call for the heavy rungs; a single ladder
  lemma then relabels them as $\partial^n\log Z$ for the textbook statement.
\item \textbf{\texttt{EventuallyEq.iteratedDeriv\_eq} is the clean tool for
  ``germ-local'' iterated derivatives.} It turns one eventual equality
  ($\deriv K =_{\text{ev}} M$) into equality of all iterated derivatives at the
  point, no induction required.
\end{itemize}

\section{Follow-ups}

\begin{itemize}[itemsep=3pt]
\item \textbf{Moment-form corollaries.} Express each $\kappa_n$ in the
  $\langle u^k\rangle_0$ (\texttt{gibbsExp}) form via moment-normalisation
  ($\kappa_2$ already done in the second-cumulant tide).
\item \textbf{Higher cumulants $\kappa_5,\dots$} if wanted: each is one more
  identical differentiation layer.
\item \textbf{Promote \texttt{amgm\_t\_abs\_x} to \texttt{lean-common}.}
\end{itemize}

\end{document}
